\documentclass[aspectratio=169]{beamer}

% Shared Beamer settings for Elements of AIML lectures (UPES).
% Keep this file free of \documentclass to allow reuse via \input.

% Allow lecture sources to use shared theme/assets without copying them into each lecture folder.
% NOTE: These paths are relative to `Unit-xx/Lecture-yy/latex/` (and `Lab/Experiment-xx/latex/`).
\makeatletter
\def\input@path{{../../../_shared/}{../../../_shared/UPES-Beamer-Theme/}}
\makeatother

\usetheme{UPES}
\setbeamertemplate{navigation symbols}{}

\usepackage{amsmath}
\usepackage{amssymb}
\usepackage{booktabs}
\usepackage{graphicx}
\usepackage{hyperref}
\usepackage{xcolor}
\usepackage{listings}

% Code listings (general-purpose).
\lstset{
  basicstyle=\ttfamily\small,
  keywordstyle=\color{blue},
  commentstyle=\color{gray},
  stringstyle=\color{teal},
  showstringspaces=false,
  columns=fullflexible,
  breaklines=true
}

\usepackage{tikz}
\usetikzlibrary{arrows.meta,positioning,shapes.geometric,calc}

% Per-slide source line (references shown on the slide itself).
\newcommand{\slidesources}[1]{%
  \vfill
  {\tiny\color{gray}\raggedright\sloppy\parbox{\linewidth}{Sources: #1}\par}%
}

% Unit-wise source strings for this subject (used by slides/notes).
% Path is relative to the lecture `latex/` folder (the compilation CWD).
% Source strings for Elements of AIML (used by slides + notes).

\newcommand{\AIMLRepoURL}{https://github.com/tali7c/Elements-of-AIML}

\newcommand{\AIMLLabSources}{%
Provost and Fawcett, \textit{Data Science for Business}.;%
McKinney, \textit{Python for Data Analysis}.;%
WEKA Documentation (Waikato Environment for Knowledge Analysis), accessed 2026-02-28.;%
Matplotlib and Seaborn documentation, accessed 2026-02-28.%
}

\newcommand{\AIMLUnitISources}{%
Rich and Knight, \textit{Artificial Intelligence}, McGraw Hill, 2017.;%
Russell and Norvig, \textit{Artificial Intelligence: A Modern Approach} (4e), Ch. 1--2 (Introduction, Intelligent Agents).;%
Poole and Mackworth, \textit{Artificial Intelligence: Foundations of Computational Agents} (2e), selected sections.%
}

\newcommand{\AIMLUnitIISources}{%
Russell and Norvig, \textit{Artificial Intelligence: A Modern Approach} (4e), Ch. 7--9 (Logical Agents, First-Order Logic, Inference).;%
Huth and Ryan, \textit{Logic in Computer Science} (2e), selected sections on propositional and first-order logic.;%
Genesereth and Nilsson, \textit{Logical Foundations of Artificial Intelligence}, selected chapters.%
}

\newcommand{\AIMLUnitIIISources}{%
Mueller and Massaron, \textit{Machine Learning for Dummies}, For Dummies, 2016.;%
Mitchell, \textit{Machine Learning}, selected chapters on learning basics and evaluation.;%
Scikit-learn documentation, accessed 2026-02-28.%
}

\newcommand{\AIMLUnitIVSources}{%
Mueller and Massaron, \textit{Machine Learning for Dummies}, For Dummies, 2016.;%
James, Witten, Hastie, and Tibshirani, \textit{An Introduction to Statistical Learning} (2e), selected chapters.;%
Scikit-learn documentation, accessed 2026-02-28.%
}

\newcommand{\AIMLUnitVSources}{%
NIST, \textit{AI Risk Management Framework (AI RMF 1.0)}, 2023.;%
Barocas, Hardt, and Narayanan, \textit{Fairness and Machine Learning} (online book), selected sections.;%
Knaflic, \textit{Storytelling with Data}, selected chapters on communicating results.%
}



\title[Elements of AIML]{Elements of AIML}
\subtitle{Miscellaneous -- Lecture 01: PCA, LDA, and ICA in One View}
\author{Tofik Ali}
\institute{School of Computer Science, UPES Dehradun}
\date{\today}

\graphicspath{{../images/}{../../../_shared/}{../../../_shared/UPES-Beamer-Theme/}}

\begin{document}

\UPESCoverFrame
\setcounter{framenumber}{0}

\begin{frame}
  \titlepage
  \vspace{0.3em}
  \begin{center}
    \footnotesize Supplementary lecture: \textbf{what each method optimizes, when to use it, and how they differ}\\[0.25em]
    \footnotesize Repository: \texttt{\AIMLRepoURL}
  \end{center}
  \slidesources{\AIMLMiscSources}
\end{frame}

\begin{frame}{Agenda}
  \tableofcontents
\end{frame}

\section{Context}

\begin{frame}{Learning Outcomes}
  By the end of this lecture, you should be able to:
  \begin{itemize}[<+->]
    \item explain the core objective of \textbf{PCA}, \textbf{LDA}, and \textbf{ICA},
    \item identify whether a method is supervised or unsupervised,
    \item compare the three methods across assumptions, outputs, and use-cases,
    \item choose a method more carefully for a new AIML problem.
  \end{itemize}
\end{frame}

\begin{frame}{Reading Order and Prerequisites}
  \begin{itemize}[<+->]
    \item First read \textbf{Unit 03 Lecture 03} for the PCA pipeline and the leakage-safe preprocessing rule.
    \item For \textbf{LDA}, basic classification language from \textbf{Unit 04} is useful because labels matter.
    \item For the eigenvalue derivation itself, use \textbf{Miscellaneous Lecture 02}.
  \end{itemize}
\end{frame}

\begin{frame}{Three Questions Before Choosing a Method}
  \begin{enumerate}[<+->]
    \item Do I have class labels?
    \item Do I want \textbf{maximum variance}, \textbf{class separation}, or \textbf{independent sources}?
    \item Is my main goal compression, classification support, visualization, or source separation?
  \end{enumerate}
\end{frame}

\section{Method Intuition}

\begin{frame}{PCA: What Does It Do?}
  \begin{block}{Principal Component Analysis}
    Find orthogonal directions that capture the \textbf{largest variance} in the data.
  \end{block}
  \begin{itemize}[<+->]
    \item Unsupervised: uses only the feature matrix.
    \item Produces new components that are linear combinations of the original features.
    \item Common goals: compression, denoising, and visualization.
  \end{itemize}
\end{frame}

\begin{frame}{LDA: What Does It Do?}
  \begin{block}{Linear Discriminant Analysis}
    Find projections that make classes far apart while keeping samples of the same class close together.
  \end{block}
  \begin{itemize}[<+->]
    \item Supervised: requires class labels.
    \item Uses between-class scatter and within-class scatter.
    \item Common goal: prepare features for classification.
  \end{itemize}
\end{frame}

\begin{frame}{ICA: What Does It Do?}
  \begin{block}{Independent Component Analysis}
    Recover latent components that are as \textbf{statistically independent} as possible.
  \end{block}
  \begin{itemize}[<+->]
    \item Unsupervised, but not a variance-maximization method.
    \item Often motivated by source separation: mixed signals $\rightarrow$ original sources.
    \item Common goals: artifact removal, blind source separation, latent-signal recovery.
  \end{itemize}
\end{frame}

\section{Mathematics}

\begin{frame}{Core Objective Functions}
  \begin{itemize}[<+->]
    \item \textbf{PCA:} maximize variance of projections.
    \[
      \max_{\|v\|=1} v^\top C v
    \]
    \item \textbf{LDA:} maximize class separation relative to class spread.
    \[
      J(W) = \frac{|W^\top S_B W|}{|W^\top S_W W|}
    \]
    \item \textbf{ICA:} maximize independence, usually through non-Gaussianity after centering and whitening.
  \end{itemize}
\end{frame}

\begin{frame}{What Matrix Problem Appears?}
  \begin{itemize}[<+->]
    \item \textbf{PCA:} eigenvectors of the covariance matrix $C$.
    \item \textbf{LDA:} generalized eigenvectors of $S_W^{-1} S_B$.
    \item \textbf{ICA:} often starts with whitening, but the final optimization is not just a covariance-eigenvalue problem.
  \end{itemize}
  \vspace{0.3em}
  \small This is a key difference: PCA and LDA are directly tied to eigen-analysis; ICA goes beyond second-order statistics.
\end{frame}

\begin{frame}{Dimensionality Limits}
  \centering
  \begin{tabular}{lll}
    \toprule
    \textbf{Method} & \textbf{Label use} & \textbf{Dimension limit} \\
    \midrule
    PCA & No & choose any $k \le d$ \\
    LDA & Yes & at most $C-1$ for $C$ classes \\
    ICA & No & usually same order as observed signals/features \\
    \bottomrule
  \end{tabular}
\end{frame}

\section{Detailed Comparison}

\begin{frame}{Detailed Comparison I}
  \centering
  \begin{tabular}{p{2.1cm}p{2.2cm}p{3.1cm}p{3.5cm}}
    \toprule
    \textbf{Method} & \textbf{Supervised?} & \textbf{Optimizes} & \textbf{Best used for} \\
    \midrule
    PCA & No & maximum variance & compression, visualization, denoising \\
    LDA & Yes & class separation & classification-oriented projection \\
    ICA & No & independence & source separation, latent-signal extraction \\
    \bottomrule
  \end{tabular}
\end{frame}

\begin{frame}{Detailed Comparison II}
  \centering
  \begin{tabular}{p{2.1cm}p{2.6cm}p{2.7cm}p{2.8cm}}
    \toprule
    \textbf{Method} & \textbf{Important assumption} & \textbf{Component relation} & \textbf{Typical risk} \\
    \midrule
    PCA & high variance is informative & orthogonal & loses label information \\
    LDA & class labels are meaningful & not just variance-based & weak when class overlap is high \\
    ICA & latent sources are independent and often non-Gaussian & independent, not merely orthogonal & unstable if assumptions fail \\
    \bottomrule
  \end{tabular}
\end{frame}

\begin{frame}{What Students Often Confuse}
  \begin{itemize}[<+->]
    \item PCA is \textbf{not} a classification method.
    \item LDA is \textbf{not} just ``supervised PCA''; the optimization target is different.
    \item ICA is \textbf{not} the same as PCA with more components; it cares about independence, not variance.
    \item Orthogonality and independence are different ideas.
  \end{itemize}
\end{frame}

\begin{frame}{Which Method Should I Choose?}
  \begin{itemize}[<+->]
    \item Need 2D visualization or compression before modeling? \textbf{Start with PCA.}
    \item Need a low-dimensional space that separates known classes? \textbf{Consider LDA.}
    \item Need to separate mixed signals or hidden sources? \textbf{Consider ICA.}
    \item If labels are unavailable, LDA is not applicable.
  \end{itemize}
\end{frame}

\section{Summary}

\begin{frame}{Summary}
  \begin{itemize}
    \item PCA maximizes variance and is best for compression or visualization.
    \item LDA uses labels and is best when class separation matters.
    \item ICA seeks independent latent sources and is common in signal-style problems.
    \item The right method depends on the \textbf{goal}, not just on the number of dimensions.
  \end{itemize}
  \slidesources{\AIMLMiscSources}
\end{frame}

\begin{frame}{Exit Questions}
  \begin{enumerate}
    \item Why is PCA unsupervised while LDA is supervised?
    \item Why can LDA return at most $C-1$ components?
    \item What is the difference between orthogonality and independence?
    \item Give one realistic use-case for ICA.
  \end{enumerate}
\end{frame}

\end{document}
